\section{Hardness for Pairs at Distance Two}
\label{sec:distance-two}

For a graph of $\DC$, \cref{cor:class-pairs} leaves the pairs with $d(u, v) = 2$ and $c(u, v) \geq 3$.
By \cref{prop:easy-pairs}(iv), such a pair has $\cp(u, v) \in \{c(u, v), c(u, v) + 1\}$.
In general graphs, deciding whether the smaller value is attained is NP-complete (\cref{thm:hardness-distance-two}).

\begin{proposition}
\label{prop:open-pairs}
Let $G$ be a connected graph and let $u, v$ be two distinct vertices of $G$.
Then $\cp(u, v) = c(u, v)$ if and only if some minimum $u$--$v$ cut contains a $u$--$v$ path.
Such a path has an odd number of edges, at most $c(u, v)$ of them, and at least three edges if $d(u, v) = 2$.
If the cut is $\delta(A)$, the vertices of the path lie alternately in $A$ and outside $A$.
A $u$--$v$ path $P$ is contained in some minimum $u$--$v$ cut if and only if $c_{G \setminus P}(u, v) = c(u, v) - |P|$.
\end{proposition}

\begin{proof}
A minimum $u$--$v$ cut that contains a $u$--$v$ path is a cut-path with $c(u, v)$ edges, and $\cp(u, v) \geq c(u, v)$ by \cref{lem:basic-bounds}.
Conversely, a cut-path with $c(u, v)$ edges contains a $u$--$v$ cut, which has at least $c(u, v)$ edges.
The cut-path is therefore equal to this cut, so it is a minimum $u$--$v$ cut, and it contains a $u$--$v$ path.
A minimum $u$--$v$ cut is $\delta(A)$ for a set $A$ with $u \in A$ and $v \notin A$, and a $u$--$v$ path contained in $\delta(A)$ has an odd number of edges by \cref{lem:minimum-cuts}.
Its edges are distinct edges of $\delta(A)$, so it has at most $c(u, v)$ edges.
If $d(u, v) = 2$, the path has at least two edges, hence at least three.
Every edge of the path has exactly one end in $A$, so the vertices of the path lie alternately in $A$ and outside $A$.

For the last statement, let $P$ be a $u$--$v$ path.
If $M$ is a $u$--$v$ cut of $G \setminus P$, then $M \cup P$ is a $u$--$v$ cut of $G$ with $|M| + |P|$ edges, so $c_{G \setminus P}(u, v) \geq c(u, v) - |P|$.
If $P$ is contained in a minimum $u$--$v$ cut $M'$ of $G$, then $M' \setminus P$ is a $u$--$v$ cut of $G \setminus P$ with $c(u, v) - |P|$ edges, and equality holds.
Conversely, suppose that equality holds, and let $M$ be a minimum $u$--$v$ cut of $G \setminus P$.
Then $M \cup P$ is a $u$--$v$ cut of $G$ with $c(u, v)$ edges, that is, a minimum $u$--$v$ cut of $G$ that contains $P$.
\end{proof}

In the graph $G_b$ of \cref{ex:counterexamples}, the minimum $u$--$v$ cut $\delta(\{u, s, b\}) = \{ua, ab, bv\}$ is the path $u, a, b, v$.
In a graph of diameter two, $\cp(u, v) = c(u, v) + 1$ for every pair with $d(u, v) = 2$~\cite[Theorem~15]{Michel2025MinCutPath}, so by \cref{prop:open-pairs} no minimum $u$--$v$ cut contains a $u$--$v$ path.

For the proof of \cref{thm:hardness-distance-two}, we assign to a formula in conjunctive normal form a graph whose minimum $u$--$v$ cuts correspond to the truth assignments of the formula (\cref{fig:hardness-construction}).
A minimum $u$--$v$ cut of this graph contains a $u$--$v$ path if and only if its truth assignment satisfies the formula.

\begin{definition}[The graph $G_\varphi$]
\label{def:formula-graph}
Let $\varphi$ be a formula in conjunctive normal form with the clauses $C_1, \dots, C_m$, where $m \geq 1$, over the variables $x_1, \dots, x_n$.
A clause is a nonempty list of literals, not necessarily distinct, a literal is a variable or a negated variable, and every variable occurs in at least one clause.
A \emph{truth assignment} sets every variable to true or to false, and it \emph{satisfies} $\varphi$ if every clause contains a literal that is true under it.
Let $t$ be the total number of literals in the $m$ clauses, counted with repetitions, and let $N = m + 1 + 2t$.
The vertex set of the graph $G_\varphi$ consists of two disjoint sets $K^u$ and $K^v$ with $N + 2$ vertices each and of one \emph{occurrence vertex} $o_{k,j}$ for the $j$-th literal of the clause $C_k$, for all $k$ and $j$.
Any two vertices of $K^u$ are adjacent, and so are any two vertices of $K^v$, and we call $K^u$ and $K^v$ the \emph{cliques}.
The clique $K^u$ contains the \emph{level vertices} $u_0, u_1, \dots, u_m$ and two \emph{anchors} $w_1, w_2$, and the clique $K^v$ contains the level vertices $v_0, v_1, \dots, v_m$ and two anchors $z_1, z_2$.
The vertices $u_k$ and $v_k$ form the \emph{level} $k$, and the other $2t$ vertices of each clique are its \emph{filler vertices}.
The \emph{rungs} $u_k v_k$, for $0 \leq k \leq m$, are the only edges between $K^u$ and $K^v$.
If the $j$-th literal of $C_k$ is a variable, then $o_{k,j}$ is a \emph{positive} occurrence vertex and is adjacent to $v_{k-1}$, $v_k$, $w_1$ and $w_2$.
If this literal is a negated variable, then $o_{k,j}$ is a \emph{negative} occurrence vertex and is adjacent to $u_{k-1}$, $u_k$, $z_1$ and $z_2$.
For every variable $x_i$, let $O_i$ be the set of the occurrence vertices of the literals $x_i$ and $\neg x_i$.
The vertices of $O_i$ are joined, in an arbitrary order, into a path by $|O_i| - 1$ \emph{synchronization edges}.
The graph $G_\varphi$ has no further edges, and we put $u = u_0$ and $v = v_m$.
For a truth assignment $\tau$, the set $A_\tau$ is the union of $K^u$ and of all sets $O_i$ such that $\tau$ sets $x_i$ to true.
\end{definition}

\begin{figure}[tb]
    \centering
    \includegraphics{hardness-construction.pdf}
    \caption{The graph $G_\varphi$ of \cref{ex:formula-graph} without the edges inside the cliques; the path in $\delta(A_\tau)$ is red}
    \label{fig:hardness-construction}
\end{figure}

Every occurrence vertex has exactly two neighbors in each clique, two level vertices in one clique and the two anchors of the other, and a filler vertex has no neighbor outside its clique.
The graph $G_\varphi$ is connected, because the rungs join the two cliques and every occurrence vertex has a neighbor in a clique.
It has $2(N + 2) + t$ vertices and $(N + 1)(N + 2) + (m + 1) + 4t + (t - n)$ edges, where the four terms count the edges inside the cliques, the rungs, the edges between occurrence vertices and cliques, and the synchronization edges.
In \cref{lem:cliques,lem:assignment-cuts,lem:satisfying}, the graph $G = (V, E)$ is the graph $G_\varphi$ of a formula $\varphi$ as in \cref{def:formula-graph}.

\begin{lemma}
\label{lem:cliques}
Let $A \subseteq V$ be a set with $|\delta(A)| \leq N$.
Then each of the cliques $K^u$ and $K^v$ is contained in $A$ or disjoint from $A$.
\end{lemma}

\begin{proof}
Suppose that $A$ contains exactly $a$ vertices of one of the two cliques, where $1 \leq a \leq N + 1$.
Each of these $a$ vertices is adjacent to each of the other $N + 2 - a$ vertices of the clique, and these $a(N + 2 - a)$ edges belong to $\delta(A)$.
Since $a(N + 2 - a) = N + 1 + (a - 1)(N + 1 - a) \geq N + 1$, this contradicts $|\delta(A)| \leq N$.
\end{proof}

\begin{lemma}
\label{lem:assignment-cuts}
(i) Let $A$ be a set with $K^u \subseteq A \subseteq V \setminus K^v$.
Then $\delta(A)$ consists of the $m + 1$ rungs, of two edges at every occurrence vertex and of the synchronization edges with exactly one end in $A$.
The two edges at an occurrence vertex $o$ are its edges to $K^v$ if $o \in A$, and its edges to $K^u$ otherwise.
In particular, $|\delta(A)| \geq N$.
(ii) Every vertex $y$ of $K^u$ and every vertex $y'$ of $K^v$ satisfy $c(y, y') = N$.
In particular, $c(u, v) = N$.
(iii) The minimum $u$--$v$ cuts are the sets $\delta(A_\tau)$ for the $2^n$ truth assignments $\tau$, and different truth assignments give different cuts.
\end{lemma}

\begin{proof}
(i) An edge inside a clique has both ends in $A$ or both ends outside $A$.
A rung has one end in $K^u \subseteq A$ and the other end in $K^v$, outside $A$.
An occurrence vertex $o$ has two neighbors in $K^u$ and two in $K^v$.
Its edges to $K^v$ belong to $\delta(A)$ if $o \in A$, and its edges to $K^u$ otherwise.
All remaining edges are synchronization edges.
The rungs and the edges at the $t$ occurrence vertices are $m + 1 + 2t = N$ edges of $\delta(A)$.

(ii) No synchronization edge has an end in $K^u$, so $\delta(K^u)$ is a $y$--$y'$ cut with $N$ edges by (i).
By \cref{lem:minimum-cuts} for the pair $y, y'$, a minimum $y$--$y'$ cut is $\delta(A)$ for a set $A$ with $y \in A$ and $y' \notin A$.
Since $|\delta(A)| \leq N$, \cref{lem:cliques} gives $K^u \subseteq A$ and $A \cap K^v = \emptyset$, and then $|\delta(A)| \geq N$ by (i).

(iii) By \cref{lem:minimum-cuts}, a minimum $u$--$v$ cut is $\delta(A)$ for a set $A$ with $u \in A$ and $v \notin A$, and $|\delta(A)| = N$ by (ii).
By \cref{lem:cliques} and (i), the minimum $u$--$v$ cuts are therefore the sets $\delta(A)$ with $K^u \subseteq A \subseteq V \setminus K^v$ such that no synchronization edge has exactly one end in $A$.
Let $A$ be such a set.
The synchronization edges join the vertices of $O_i$ into a path, so $O_i$ is contained in $A$ or disjoint from $A$, for every $i$.
Every occurrence vertex lies in one of the sets $O_i$, so $A = A_\tau$ for the truth assignment $\tau$ that sets $x_i$ to true if and only if $O_i \subseteq A$.
Conversely, both ends of a synchronization edge lie in the same set $O_i$, so no synchronization edge has exactly one end in $A_\tau$.
The sets $O_i$ are nonempty, because every variable occurs in $\varphi$.
Different truth assignments $\tau$ therefore give different sets $A_\tau$, and different cuts $\delta(A_\tau)$ by \cref{lem:minimum-cuts}.
\end{proof}

\begin{lemma}
\label{lem:satisfying}
Let $\tau$ be a truth assignment.
The cut $\delta(A_\tau)$ contains a $u$--$v$ path if and only if $\tau$ satisfies $\varphi$.
Every $u$--$v$ path contained in $\delta(A_\tau)$ has at least $2m + 1$ edges.
\end{lemma}

\begin{proof}
An occurrence vertex is \emph{true} if its literal is true under $\tau$, and \emph{false} otherwise.
A positive occurrence vertex of a variable $x_i$ lies in $A_\tau$ if and only if $\tau$ sets $x_i$ to true, that is, if and only if it is true.
A negative occurrence vertex lies in $A_\tau$ if and only if it is false.
By \cref{lem:assignment-cuts}, the cut $\delta(A_\tau)$ has $N$ edges, so it contains no synchronization edge and consists of the rungs and of the following two edges at every occurrence vertex $o$ of every clause $C_k$.
If $o$ is true, the two edges join $o$ to $v_{k-1}$ and $v_k$ or to $u_{k-1}$ and $u_k$, that is, to a vertex of level $k - 1$ and to a vertex of level $k$.
If $o$ is false, they join $o$ to two anchors.

For $1 \leq k \leq m$, let $B_k$ be the set of the vertices of the levels $0, \dots, k - 1$ and of the true occurrence vertices of the clauses $C_1, \dots, C_{k-1}$.
Then $u \in B_k$ and $v \notin B_k$.
Both ends of a rung lie in the same level.
An edge of $\delta(A_\tau)$ at a false occurrence vertex joins it to an anchor, and neither of them lies in $B_k$.
An edge of $\delta(A_\tau)$ at a true occurrence vertex $o$ of a clause $C_j$ joins $o$ to a vertex of level $j - 1$ or $j$.
Both ends of this edge lie in $B_k$ if $j < k$, and neither does if $j > k$ or if $j = k$ and the edge leads to level $k$.
Hence the edges of $\delta(A_\tau)$ with exactly one end in $B_k$ are the edges that join a true occurrence vertex of $C_k$ to a vertex of level $k - 1$.

Suppose that $\tau$ does not satisfy a clause $C_k$.
Then $C_k$ has no true occurrence vertex, so no edge of $\delta(A_\tau)$ has exactly one end in $B_k$, and $\delta(A_\tau)$ contains no $u$--$v$ path.
Suppose now that $\tau$ satisfies $\varphi$.
For $1 \leq k \leq m$, the clause $C_k$ has a true occurrence vertex, which is joined by edges of $\delta(A_\tau)$ to a vertex of level $k - 1$ and to a vertex of level $k$.
The two vertices of each level are joined by a rung.
By induction on $k$, the graph $(V, \delta(A_\tau))$ contains a path from $u = u_0$ to $u_k$ and to $v_k$ for $0 \leq k \leq m$, in particular a $u$--$v$ path.

Finally, let $P$ be a $u$--$v$ path contained in $\delta(A_\tau)$, and let $1 \leq k \leq m$.
The path $P$ starts in $B_k$ and ends outside $B_k$, so it contains an edge that joins a vertex of level $k - 1$ to a true occurrence vertex $o$ of $C_k$.
Since $o \notin \{u, v\}$, the path $P$ contains two edges at $o$.
Every edge of $\delta(A_\tau)$ has at most one end that is an occurrence vertex, and the occurrence vertices of different clauses are distinct, so the $2m$ edges obtained for $k = 1, \dots, m$ are distinct.
Hence $|P| \geq 2m$, and $|P| \geq 2m + 1$ because $|P|$ is odd (\cref{prop:open-pairs}).
\end{proof}

\begin{theorem}
\label{thm:hardness-distance-two}
It is NP-complete to decide, for a connected graph $G$ and two vertices $u, v$ of $G$ with $d(u, v) = 2$, whether $\cp(u, v) = c(u, v)$.
The problem remains NP-complete for graphs of diameter three.
\end{theorem}

\begin{proof}
By \cref{prop:open-pairs}, $\cp(u, v) = c(u, v)$ if and only if some $u$--$v$ path $P$ satisfies $c_{G \setminus P}(u, v) = c(u, v) - |P|$.
For a given path $P$, this equation is verified in polynomial time by two computations of a maximum flow~\cite[Sections~10.3 and~10.4]{Erickson2019Algorithms}, so the problem belongs to NP.

The satisfiability problem for formulas in conjunctive normal form with at most three literals per clause is NP-complete~\cite{Cook1971Complexity,Karp1972Reducibility}.
Let $\varphi$ be such a formula.
A formula without clauses is satisfiable and a formula with an empty clause is not, so we may assume that $\varphi$ has $m \geq 1$ clauses, all of them nonempty, and that $x_1, \dots, x_n$ are the variables that occur in $\varphi$.
By \cref{lem:assignment-cuts}(iii) and \cref{lem:satisfying}, some minimum $u$--$v$ cut of $G_\varphi$ contains a $u$--$v$ path if and only if some truth assignment satisfies $\varphi$.
By \cref{prop:open-pairs}, the pair $u, v$ of $G_\varphi$ therefore satisfies $\cp(u, v) = c(u, v)$ if and only if $\varphi$ is satisfiable.
The graph $G_\varphi$ has $2(N + 2) + t = 2m + 5t + 6$ vertices and is constructed in polynomial time.

In $G_\varphi$, the vertices $u = u_0$ and $v = v_m$ are not adjacent, because the rungs are the only edges between the cliques and $m \geq 1$.
The path $u_0, u_m, v_m$ shows that $d(u, v) = 2$.
Two vertices of the same clique are adjacent.
A vertex $y$ of $K^u$ and a vertex $y'$ of $K^v$ satisfy $d(y, y') \leq d(y, u_0) + 1 + d(v_0, y') \leq 3$.
An occurrence vertex has a neighbor in each clique, so it is at distance at most two from every vertex of a clique.
Two occurrence vertices have neighbors in $K^u$, which are equal or adjacent, so they are at distance at most three.
A filler vertex of $K^u$ and a filler vertex of $K^v$ are not adjacent and have no common neighbor, so they are at distance three.
Hence the largest distance between two vertices of $G_\varphi$ is three, that is, $G_\varphi$ has diameter three.
\end{proof}

A pair with $d(u, v) = 1$ always satisfies $\cp(u, v) = c(u, v)$ (\cref{prop:easy-pairs}(i)), and in a graph of diameter at most two no pair with $d(u, v) = 2$ does~\cite[Theorem~15]{Michel2025MinCutPath}.
Hence neither the distance two of the pair nor the diameter three of the graph can be replaced by a smaller value in \cref{thm:hardness-distance-two}.

\begin{corollary}
\label{cor:distance-two}
(i) Computing $\cp(u, v)$ is NP-hard for pairs with $d(u, v) = 2$, also in graphs of diameter three.
(ii) Deciding whether $\cp(u, v) = c(u, v) + d(u, v) - 1$ is coNP-complete, also for pairs with $d(u, v) = 2$ in graphs of diameter three.
\end{corollary}

\begin{proof}
(i) The value $c(u, v)$ is computed in polynomial time~\cite[Sections~10.3 and~10.4]{Erickson2019Algorithms}, so an algorithm that computes $\cp(u, v)$ decides whether $\cp(u, v) = c(u, v)$.

(ii) By \cref{lem:basic-bounds}, the equality fails if and only if some cut-path between $u$ and $v$ has at most $c(u, v) + d(u, v) - 2$ edges.
A set $S$ of edges is a cut-path if and only if the graph $(V, S)$ contains a $u$--$v$ path and the graph $G \setminus S$ contains none, which two searches decide.
The problem therefore belongs to coNP.
If $d(u, v) = 2$, then $\cp(u, v) \in \{c(u, v), c(u, v) + 1\}$ by \cref{prop:easy-pairs}(iv), so the equality holds if and only if $\cp(u, v) \neq c(u, v)$.
Deciding whether $\cp(u, v) \neq c(u, v)$ is coNP-hard by \cref{thm:hardness-distance-two}, also in graphs of diameter three.
\end{proof}

For a pair with $d(u, v) = 2$, a shortest $u$--$v$ path shares exactly one edge with every minimum $u$--$v$ cut (\cref{lem:minimum-cuts}).
Statement (ii) then asks whether the union of a minimum $u$--$v$ cut and a shortest $u$--$v$ path, which has $c(u, v) + 1$ edges, is a minimum cut-path.

\begin{remark}[The graphs $G_\varphi$]
\label{rem:formula-graph}
Let $\varphi$ be a satisfiable formula with $m$ clauses.
By \cref{lem:assignment-cuts}(iii) and \cref{lem:satisfying}, every $u$--$v$ path contained in a minimum $u$--$v$ cut of $G_\varphi$ has at least $2m + 1$ edges.
Satisfiable formulas exist for every $m$, so no constant $\ell$ has the property that every pair with $d(u, v) = 2$ and $\cp(u, v) = c(u, v)$ has a minimum $u$--$v$ cut that contains a $u$--$v$ path with at most $\ell$ edges.
No graph $G_\varphi$ belongs to $\DC$.
A filler vertex of $K^u$ and a filler vertex of $K^v$ are at distance three by the proof of \cref{thm:hardness-distance-two}, and their cut-value is $N \geq 4$ by \cref{lem:assignment-cuts}(ii).
\end{remark}

Whether $\cp(u, v) = d(u, v)$ is decided in polynomial time for pairs with bounded $d(u, v)$ and is NP-complete for pairs with $c(u, v) = 2$.
For the question whether $\cp(u, v) = c(u, v)$, the distance and the cut-value exchange their roles (\cref{tab:bounded}).

\begin{proposition}
\label{prop:bounded}
Let $k$ be a fixed positive integer.
(i) There is a polynomial-time algorithm that decides, for a connected graph $G$ and two vertices $u, v$ of $G$ with $d(u, v) \leq k$, whether $\cp(u, v) = d(u, v)$.
(ii) There is a polynomial-time algorithm that decides, for a connected graph $G$ and two vertices $u, v$ of $G$ with $c(u, v) \leq k$, whether $\cp(u, v) = c(u, v)$.
\end{proposition}

\begin{proof}
A $u$--$v$ path with at most $k$ edges is determined by the sequence of its at most $k - 1$ vertices other than $u$ and $v$.
Hence $G = (V, E)$ contains at most $k|V|^{k-1}$ such paths, and they are listed in polynomial time.

(i) By \cref{lem:basic-bounds}, $\cp(u, v) = d(u, v)$ if and only if some shortest $u$--$v$ path $P$ is separating.
Every shortest $u$--$v$ path has $d(u, v) \leq k$ edges, and one search in $G \setminus P$ decides whether $P$ is separating.

(ii) By \cref{prop:open-pairs}, $\cp(u, v) = c(u, v)$ if and only if some $u$--$v$ path $P$ with at most $c(u, v) \leq k$ edges satisfies $c_{G \setminus P}(u, v) = c(u, v) - |P|$.
For each path $P$, this equation is tested by one computation of a maximum flow in $G \setminus P$.
\end{proof}

\begin{table}[tb]
\caption{Deciding $\cp(u, v) = d(u, v)$ and $\cp(u, v) = c(u, v)$ for a fixed integer $k \geq 2$}\label{tab:bounded}
\begin{tabular}{@{}lll@{}}
\toprule
Pairs & $\cp(u, v) = d(u, v)$ & $\cp(u, v) = c(u, v)$ \\
\midrule
Bounded cut-value, $c(u, v) \leq k$ & NP-complete (\cref{thm:hardness}) & polynomial (\cref{prop:bounded}(ii)) \\
Bounded distance, $d(u, v) \leq k$ & polynomial (\cref{prop:bounded}(i)) & NP-complete (\cref{thm:hardness-distance-two}) \\
\botrule
\end{tabular}
\end{table}

In general graphs, computing $\cp(u, v)$ is thus NP-hard for pairs with $c(u, v) = 2$ and for pairs with $d(u, v) = 2$.
By \cref{def:class}, every pair of a graph of $\DC$ is of one of these two kinds or has $d(u, v) = 1$ or $c(u, v) = 1$.

\begin{example}
\label{ex:formula-graph}
Let $\varphi = (x_1 \vee \neg x_2) \wedge x_2$ (\cref{fig:hardness-construction}).
Then $m = 2$, $t = 3$ and $N = 9$, each clique has $11$ vertices, six of which are filler vertices, and $G_\varphi$ has $25$ vertices and $110 + 3 + 12 + 1 = 126$ edges.
The occurrence vertices are $o_{1,1}$ (positive, for $x_1$), $o_{1,2}$ (negative, for $x_2$) and $o_{2,1}$ (positive, for $x_2$).
The only synchronization edge joins the two vertices of $O_2 = \{o_{1,2}, o_{2,1}\}$.
By \cref{lem:assignment-cuts}, $c(u, v) = 9$ and the pair $u, v$ has four minimum cuts.
Only the truth assignment $\tau$ that sets $x_1$ and $x_2$ to true satisfies $\varphi$.
The cut $\delta(A_\tau)$ consists of the three rungs and of the edges from $o_{1,1}$ to $v_0$ and $v_1$, from $o_{2,1}$ to $v_1$ and $v_2$, and from $o_{1,2}$ to $z_1$ and $z_2$.
It contains the path $P = u_0, v_0, o_{1,1}, v_1, o_{2,1}, v_2$ with $2m + 1 = 5$ edges, and $c_{G_\varphi \setminus P}(u, v) = 4 = 9 - 5$.
Hence $\cp(u, v) = c(u, v) = 9$, whereas $d(u, v) = 2$ and $c(u, v) + d(u, v) - 1 = 10$.
\end{example}

For graphs of $\DC$ and pairs with $d(u, v) = 2$ and $c(u, v) \geq 3$, we do not know whether $\cp(u, v)$ can be computed in polynomial time or whether computing it is NP-hard.
\Cref{thm:hardness-distance-two} does not answer this question, because no graph $G_\varphi$ belongs to $\DC$ (\cref{rem:formula-graph}).
